Napjaink felgyorsult világában kevesen használják a tömegközlekedést, mert az nem elég kényelmes, nem elég gyors vagy éppen nem úgy működik, ahogy azt elvárnánk tőle. A gyakori problémák közé tartozik, hogy a buszok nem tartják a menetrendet vagy a jegyek megvásárlása nehézkes. Sok helyen a menetrend nincs kifüggesztve a megállókban és az interneten sem elérhető. Az is gyakran megesik hogy a buszok nagyon zsúfoltak és erről az utazni vágyók csak akkor vesznek tudomást, amikor már a megállóban vannak és nem áll rendelkezésre más utazási lehetőség. A Bus4U online platform ezen problémákra kínál megoldást.

A Bus4U egy mobilalkalmazásból, egy felhasználói weboldalból és egy adminisztrátori weboldalból tevődik össze. Főbb funkciói közé tartozik a buszjáratok keresése a kiindulópont, a célállomás és az időpont megadásával, a járatokra szóló előzetes online jegyvásárlás, illetve a járatok indulási időpontjainak megtekintése városokra és megállókra lebontva. Ezen kívül megtekinthetőek a buszmegállók a térképen, városok és útvonalak szerint szűrve, valamint valós időben követhetőek az autóbuszok pillanatnyi helyzetei is. Ezen funkciók biztosításához, illetve a jegykezeléshez, szükség van egy alkalmazásra, ami a buszokon található POS terminálon fut. Ez biztosítja a jegyek érvényesítését, illetve az autóbuszok élő információinak megosztását is a központi szerverrel.

Ezen felül a webalkalmazás biztosít egy adminisztrációs felületet is a busztársaságok számára. Itt lehetőség nyílik elsősorban a menetrendek és a megállók kezelésére, valamint az útvonalak feltöltésére. A busztársaság adminisztrátora itt tudja kezelni az alkalmazottak információit, és az autóbuszok különböző adatait is, mint az időszakos műszaki vizsga és a biztosítás. Ugyanitt található egy statisztikai felület is, ahol a busztársaságok vezetői megtekinthetik a regisztrált felhasználók számát, az eladott jegyek számát és az ezekből generált bevételt, mindezt havi, éves és mindenkori bontásban.

Továbbfejlesztési lehetőségek között szerepelnek a bérletek létrehozása és a buszok telítettségének követése valós időben. A busztársaságok szempontjából fontos lehet még a sofőrök munkaidejének követése és az automatizált könyvelés bevezetése is.

Ez a dolgozat a szoftver kliens oldali alkalmazásaira fókuszál, a két weboldalra és a mobilalkalmazásra, míg a backend működése egy másik dolgozat keretein belül van ismertetve.

\textbf{Kulcsszavak:} tömegközlekedés, buszjárat, digitalizálás